% BEGIN REV09_FORMA_PARTIDA_ENERGIA
\section{Forma partida, integralidad y energía reticular}
\label{corpus9:xii:forma-partida}

La deformación del teorema~\ref{rec:thm:red} parte de la red marcada
\(\Lambda\), de la dirección \(u=P_3P_{11}K\) y de las coordenadas
regionales previamente generadas. Denotamos por \(E\) su espacio real
de dimensión veinticuatro y por \(\langle\,\cdot\,,\,\cdot\,\rangle_B\)
la métrica positiva de los doce planos ortogonales de tipo \(A_2\).
Conservamos la normalización de \eqref{rec:eq:flujo} y
\eqref{rec:eq:red}:
\begin{equation}
 \kappa=\frac{\log\rho}{\|u\|},\qquad h_i=\kappa u_i,
 \qquad
 G_t=\bigoplus_{i=1}^{12}e^{th_i}I_{A_{2,i}},\qquad
 \Lambda_t=G_t\Lambda.
 \label{corpus9:xii:normalizacion-flujo}
\end{equation}
El operador \(G_t\) es autoadjunto para \(B\), y
\(G_t^{-1}=G_{-t}\). Si se escribe la misma colección con el parámetro
sin normalizar, \(D_s=\bigoplus_i e^{su_i}I_{A_{2,i}}\), la relación
exacta es \(G_t=D_{\kappa t}\). El parámetro de deformación y el
reloj TPK mantienen sus dominios respectivos.

\subsection{Variación aritmética de la proyección euclídea}
\label{corpus9:xii:variacion-euclidea}

La coordenatización ordenada del lema~\ref{lem:k-sector-no-nulo} determina,
por la suma finita \eqref{eq:k-proyectores-incidencia},
\begin{equation}
 u_1=-\frac{495}{4}-\frac{233}{10}\sqrt5,
 \qquad \sum_{i=1}^{12}u_i=0.
 \label{corpus9:xii:coordenada-testigo}
\end{equation}
La normalización regional cumple \(\rho=\pi_{\HMT}/\varphi_{\HMT}^2>1\):
en la realización ya establecida, \(\pi_{\HMT}>3\) y
\(\varphi_{\HMT}^2=(3+\sqrt5)/2<3\). Por tanto, \(\kappa>0\).
Estas comparaciones se efectúan sobre las salidas regionales, después
de su generación.

\begin{proposition}[Testigo local de variación de la integralidad]
\label{corpus9:xii:testigo-integralidad}
Existe \(\delta>0\) tal que, para \(0<|t|<\delta\), la red
\(\Lambda_t\) contiene un vector de norma cuadrada no entera.
En consecuencia, su métrica euclídea carece de integralidad y de una
isometría con el retículo de Leech.
\end{proposition}
\begin{proof}
El vector marcado de \eqref{exc:vector} es
\(v_\alpha=4\rho_1+\sum_{j=2}^{12}\rho_j\), con
\(\|\rho_j\|_B^2=2\) en planos mutuamente ortogonales y
\(v_\alpha^2=54\). La construcción del vecino
\eqref{exc:vecino} da \(v_\alpha/3\in\Lambda\). Su norma transportada es
\begin{equation}
 f(t)=\left\|G_t\frac{v_\alpha}{3}\right\|_B^2
 =\frac29\left(16e^{2th_1}+\sum_{j=2}^{12}e^{2th_j}\right).
 \label{corpus9:xii:norma-testigo}
\end{equation}
La suma nula y \eqref{corpus9:xii:coordenada-testigo} producen
\begin{equation}
 \begin{aligned}
 f(0)&=6,\\
 f'(0)&=\frac49\left(16h_1+\sum_{j=2}^{12}h_j\right)
       =\frac{20}{3}\kappa u_1
       =\kappa\left(-825-\frac{466}{3}\sqrt5\right)\ne0.
 \end{aligned}
 \label{corpus9:xii:derivada-testigo}
\end{equation}
La continuidad de \(f'\) proporciona un intervalo
\((-\delta,\delta)\) en el que \(f\) es estrictamente monótona.
Reduciendo \(\delta\), se obtiene además
\(11/2<f(t)<13/2\). Para \(0<|t|<\delta\), la monotonía da
\(f(t)\ne f(0)=6\); el único entero en ese intervalo es seis.
Así, \(G_t(v_\alpha/3)\) tiene norma cuadrada no entera. La
integralidad de una red exige normas cuadradas enteras, y las isometrías
conservan esas normas. La integralidad de Leech prueba la última
conclusión.
\end{proof}

La deformación conserva el covolumen y la simetría ternaria mientras
varía su estructura métrica aritmética. La reciprocidad
\(\Lambda_t^*=\Lambda_{-t}\) utiliza la autoadjunción del flujo y
la autodualidad inicial. Una permutación ambiental que intercambia
dilataciones y contracciones actúa sobre el retículo pegado cuando
preserva además el código y el vecino que lo construyen.

\subsection{Autodualidad integral para la forma de signatura partida}
\label{corpus9:xii:autodualidad-partida}

La consideración simultánea de ambas proyecciones utiliza el espacio
\(E\oplus E\), con forma bilineal
\begin{equation}
 \mathcal B((x,p),(x',p'))
 =\langle x,p'\rangle_B+\langle p,x'\rangle_B.
 \label{corpus9:xii:forma-bilineal-partida}
\end{equation}
Definimos
\begin{equation}
 \begin{aligned}
 \mathcal G_t&=G_t\oplus G_{-t},\\
 \Gamma_t&=\mathcal G_t(\Lambda\oplus\Lambda)
 =\{(G_t\lambda,G_{-t}\mu):\lambda,\mu\in\Lambda\}.
 \end{aligned}
 \label{corpus9:xii:red-partida}
\end{equation}

\begin{theorem}[Integralidad, paridad y autodualidad del retículo partido]
\label{corpus9:xii:teorema-partido}
La forma \(\mathcal B\) tiene signatura \((24,24)\).
Para todo \(t\in\mathbb R\), \(\mathcal G_t\) es una isometría
de \(\mathcal B\), y \(\Gamma_t\) es una red integral, par y
autodual respecto de esa forma.
\end{theorem}
\begin{proof}
Las coordenadas
\[
 p_L=\frac{x+p}{\sqrt2},\qquad p_R=\frac{x-p}{\sqrt2}
\]
constituyen un cambio lineal invertible, con
\(x=(p_L+p_R)/\sqrt2\) y \(p=(p_L-p_R)/\sqrt2\). Satisfacen
\[
 \mathcal B((x,p),(x,p))=\|p_L\|_B^2-\|p_R\|_B^2.
\]
La positividad de \(B\) en dimensión veinticuatro da la signatura
\((24,24)\). La autoadjunción de \(G_t\) y la identidad
\(G_tG_{-t}=I\) implican
\[
 \langle G_tx,G_{-t}p'\rangle_B
 =\langle x,G_tG_{-t}p'\rangle_B=\langle x,p'\rangle_B.
\]
El otro producto cruzado satisface la misma igualdad. Por tanto,
\(\mathcal B(\mathcal G_t\xi,\mathcal G_t\eta)
=\mathcal B(\xi,\eta)\).

En \(\Gamma_0=\Lambda\oplus\Lambda\), los productos cruzados son
enteros, porque \(\Lambda\) es integral, y
\[
 \mathcal B((\lambda,\mu),(\lambda,\mu))
 =2\langle\lambda,\mu\rangle_B\in2\mathbb Z.
\]
Esto prueba integralidad y paridad. Si \((x,p)\) pertenece al dual
de \(\Gamma_0\) respecto de \(\mathcal B\), sus emparejamientos
con \((\lambda,0)\) y \((0,\mu)\), para todo
\(\lambda,\mu\in\Lambda\), exigen respectivamente
\(p,x\in\Lambda^*=\Lambda\). Recíprocamente, esas dos pertenencias
garantizan todos los emparejamientos enteros. Así,
\(\Gamma_0^{*\mathcal B}=\Gamma_0\). La isometría
\(\mathcal G_t\) transporta las tres propiedades y da
\(\Gamma_t^{*\mathcal B}=\mathcal G_t\Gamma_0^{*\mathcal B}
=\Gamma_t\).
\end{proof}

La métrica determina el dominio de cada conclusión. Los productos
cruzados de \(\mathcal B\) conservan la integralidad de
\(\Gamma_t\), incluso para los parámetros de la
proposición~\ref{corpus9:xii:testigo-integralidad}. Los retículos
\(\Gamma_t\), considerados con esa forma partida, son isométricos
entre sí. Sus proyecciones euclídeas \(\Lambda_t\) y
\(\Lambda_{-t}\), equipadas con \(B\), varían con el parámetro.

\begin{proposition}[Intercambio de las proyecciones recíprocas]
\label{corpus9:xii:involucion-partida}
La involución \(\mathcal J(x,p)=(p,x)\) satisface
\[
 \mathcal J\Gamma_t=\Gamma_{-t},\qquad
 \mathcal J\mathcal G_t\mathcal J=\mathcal G_{-t}.
\]
En las coordenadas \((p_L,p_R)\), fija \(p_L\) e invierte
el signo de \(p_R\).
\end{proposition}
\begin{proof}
El intercambio transforma \((G_t\lambda,G_{-t}\mu)\) en
\((G_{-t}\mu,G_t\lambda)\). Al recorrer \(\lambda,\mu\in\Lambda\),
la imagen recorre \(\Gamma_{-t}\). La conjugación intercambia los
dos bloques de \(\mathcal G_t\), produciendo \(\mathcal G_{-t}\).
Finalmente, \(x+p\) se conserva y \(x-p\) cambia de signo.
\end{proof}

Esta involución actúa sobre la construcción reticular completa. Una
representación física posterior debe transportar también los operadores
y los datos de identificación propios de esa realización.

\subsection{Operador positivo y equivalencia unitaria de dominios}
\label{corpus9:xii:energia-reticular}

La energía de las dos proyecciones se define sobre la red fija por
\begin{equation}
 \mathsf E_t^{\rm lat}(\lambda,\mu)
 =\|G_t\lambda\|_B^2+\|G_{-t}\mu\|_B^2,
 \qquad \lambda,\mu\in\Lambda.
 \label{corpus9:xii:energia}
\end{equation}
Es no negativa, se anula exactamente en \((0,0)\) y satisface
\(\mathsf E_t^{\rm lat}(\lambda,\mu)
=\mathsf E_{-t}^{\rm lat}(\mu,\lambda)\). El intercambio de ambas
contribuciones expresa, en los doce planos marcados, la reciprocidad
de una suma de términos de escala opuesta. Para
\(M_h=\max_i|h_i|\), la descomposición ortogonal en planos da
\begin{equation}
 e^{-2|t|M_h}\|v\|_B^2\leq\|G_tv\|_B^2
                    \leq e^{2|t|M_h}\|v\|_B^2.
 \label{corpus9:xii:cotas-cuadraticas}
\end{equation}
En efecto, cada sumando \(\|v_i\|_B^2\) se multiplica por
\(e^{2th_i}\), comprendido entre esas dos cotas.

\begin{theorem}[Autoadjunción, resolvente compacto y dualidad unitaria]
\label{corpus9:xii:operador-positivo}
En \(\ell^2(\Lambda\oplus\Lambda)\), el operador
\[
 (H_t^{\rm lat}\psi)(\lambda,\mu)
 =\mathsf E_t^{\rm lat}(\lambda,\mu)\psi(\lambda,\mu)
\]
con dominio máximo
\begin{equation}
 \mathcal D(H_t^{\rm lat})=
 \left\{\psi\in\ell^2(\Lambda\oplus\Lambda):
 \sum_{\lambda,\mu}\bigl(\mathsf E_t^{\rm lat}(\lambda,\mu)\bigr)^2
                  |\psi(\lambda,\mu)|^2<\infty\right\}
 \label{corpus9:xii:dominio-operador}
\end{equation}
es positivo, autoadjunto y de resolvente compacto. La aplicación
\(U\psi(\lambda,\mu)=\psi(\mu,\lambda)\) es unitaria y cumple
\begin{equation}
 U\mathcal D(H_t^{\rm lat})=\mathcal D(H_{-t}^{\rm lat}),\qquad
 UH_t^{\rm lat}U^{-1}=H_{-t}^{\rm lat}.
 \label{corpus9:xii:equivalencia-unitaria}
\end{equation}
\end{theorem}
\begin{proof}
Las funciones de soporte finito pertenecen al dominio y forman un
subespacio denso. La energía es real. Si \(\psi\) pertenece al dominio
del adjunto, evaluar su funcional contra cada vector de la base
ortonormal de soporte unitario fuerza que la coordenada del adjunto
sea \(\mathsf E_t^{\rm lat}(\lambda,\mu)\psi(\lambda,\mu)\).
La pertenencia de esa sucesión a \(\ell^2\) es exactamente la condición
\eqref{corpus9:xii:dominio-operador}. Recíprocamente, dicha condición
permite el emparejamiento con toda función del dominio por
Cauchy--Schwarz. El dominio del adjunto y su acción coinciden con los
de \(H_t^{\rm lat}\), lo que prueba autoadjunción. Además,
\[
 \langle\psi,H_t^{\rm lat}\psi\rangle
 =\sum_{\lambda,\mu}\mathsf E_t^{\rm lat}(\lambda,\mu)
                         |\psi(\lambda,\mu)|^2\geq0.
\]
Por \eqref{corpus9:xii:cotas-cuadraticas},
\[
 \mathsf E_t^{\rm lat}(\lambda,\mu)
 \geq e^{-2|t|M_h}
       \bigl(\|\lambda\|_B^2+\|\mu\|_B^2\bigr).
\]
La discreción de \(\Lambda\oplus\Lambda\) implica que cada subnivel
de la energía es finito. El resolvente \((H_t^{\rm lat}+I)^{-1}\)
es diagonal, con coeficientes \((1+\mathsf E_t^{\rm lat})^{-1}\).
Su restricción al subnivel \(\mathsf E_t^{\rm lat}\leq R\) tiene
rango finito y el resto tiene norma a lo sumo \((1+R)^{-1}\).
El resolvente es, por tanto, límite en norma de operadores de rango
finito y es compacto.

El intercambio de índices conserva la norma de \(\ell^2\) y satisface
\(U^{-1}=U\). Al cambiar \((\lambda,\mu)\) por \((\mu,\lambda)\)
en la suma que define el dominio y utilizar la igualdad de energías,
se obtiene la primera identidad de
\eqref{corpus9:xii:equivalencia-unitaria}. Sobre ese dominio,
\[
 (UH_t^{\rm lat}U^{-1}\psi)(\lambda,\mu)
 =\mathsf E_t^{\rm lat}(\mu,\lambda)\psi(\lambda,\mu)
 =\mathsf E_{-t}^{\rm lat}(\lambda,\mu)\psi(\lambda,\mu),
\]
que prueba la identidad de operadores con sus dominios.
\end{proof}

Las tres conservaciones quedan así tipadas. La métrica euclídea
conserva covolumen y simetría ternaria y relaciona los parámetros
opuestos mediante la red dual. La forma partida conserva integralidad,
paridad y autodualidad de la red duplicada. La realización operatoria
conserva la energía mediante una equivalencia unitaria que transporta
su dominio. La lectura theta del teorema~\ref{rec:thm:theta} expresa
la reciprocidad en la suma convergente de longitudes, y la continuación
Mellin conserva su estructura polar por el
teorema~\ref{rec:thm:mellin}.
% END REV09_FORMA_PARTIDA_ENERGIA
